\section{The Klein quartic in an eight-by-eight slice}
\label{sec:klein}
This section identifies the permanent level conditions on a conjugate pair of three-dimensional representations of $\operatorname{PSL}_2(\mathbb F_7)$. Klein's invariant ring supplies the geometry; the permanent calculations determine which of its invariants occur.

Put $s=\sqrt{-7}$ and $\alpha=(-1+s)/2$. Index rows and columns by
$P=\mathbb F_7\cup\{\infty\}$, in that order, and let
$\rho(g)T=P_gTP_g^{\mathsf T}$. Let $u:z\mapsto z+1$ and
$v:z\mapsto-1/z$, and let $C$ be the conjugacy class of $u$. Define
\[
 D=\sum_{g\in C}\rho(g)-\sum_{g\in C}\rho(g^{-1}),\qquad
 W=\ker(D-8sI).
\]
The character table in Elkies~\cite[\S1.1, p.~53, (1.1)]{Elkies1999}
shows that $W$ is one of the two non-real irreducible summands of dimension three in $\operatorname{Mat}_8(\mathbb C)$.

Here are coordinates over $\mathbb Q(s)$, so the normalization of every coefficient below is specified. In permutation-tuple notation put
$a=(1,0,4,7,2,6,5,3)$, let $h:z\mapsto4z$, and set
\[
 B_0=\frac18(I+\rho(a))(D^2+8sD)E_{02},\quad
 B_1=\rho(h)B_0,\quad B_2=\rho(h^2)B_0.
\]
Write $B(x)=x_0B_0+x_1B_1+x_2B_2$ and
\begin{align*}
 F(x)&=x_0^4+x_1^4+x_2^4
       +3\alpha(x_0^2x_1^2+x_0^2x_2^2+x_1^2x_2^2),\\
 H(x)&=\det\left(\frac{\partial^2F}{\partial x_i\partial x_j}\right)_{0\leq i,j\leq2},\\
 c&=2688(31+3s),\qquad d=-\frac{832}{27}(87+91s).
\end{align*}
A bar on a polynomial or matrix conjugates its coefficients, leaving its variables fixed.

\begin{theorem}\label{thm:klein-pure}
The matrices $B_0,B_1,B_2$ are a basis of $W$, and
\[
 S_1|_W=S_2|_W=S_3|_W=S_5|_W=S_7|_W=0,
 \qquad S_4(B(x))=cF(x),\qquad S_6(B(x))=dH(x),
\]
\[
 S_8|_W=\frac5{1764}(S_4|_W)^2.
\]
The projective level locus in $\mathbb P(W)$ consists of exactly the twenty-four ordinary flexes of Klein's quartic $F=0$. In particular no two distinct level lines in $W$ span a level plane. The same assertions hold on $\overline W$.
\end{theorem}

\begin{proof}
The conjugacy classes are $1A$, $2A$, $3A$, $4A$, $7A$ and $7B$.
The permutation character on $P$ has values $(8,0,2,0,1,1)$ in this order. Its square has inner product
$(192-24)/168=1$ with each three-dimensional character. The central operator
$D$ acts on these summands by $8s$ and $-8s$, respectively, and by zero on all other irreducibles. Thus $W$ has dimension three. The symmetric square of the permutation representation has character $(36,4,3,0,1,1)$; its inner product with either three-dimensional character is
$(108-84-24)/168=0$. Hence the matrices in $W$ are skew-symmetric. Their row and column sums vanish because those sums define equivariant maps to the permutation representation, which has no three-dimensional constituent.

Direct evaluation of the displayed class-sum expression gives
\[
 \sum_{r,j}(B_i)_{rj}(B_k)_{rj}=0,
 \qquad
 \sum_{r,j}(B_i)_{rj}(\overline B_k)_{rj}=1792\delta_{ik}.
\]
In particular the three matrices are independent. Their generator matrices are
\[
 \rho(u)|_W=
 \begin{pmatrix}
 \alpha/2&\alpha/2&0\\
 1/2&-1/2&\overline\alpha/2\\
 1/2&-1/2&-\overline\alpha/2
 \end{pmatrix},\qquad
 \rho(v)|_W=
 \begin{pmatrix}0&0&-1\\0&-1&0\\-1&0&0\end{pmatrix}.
\]
The quartic $F$ is invariant under these matrices. Klein's invariant ring,
as described by Elkies~\cite[\S1.2, pp.~55--58, (1.9)--(1.18)]{Elkies1999},
has generators of degrees $4,6,14,21$. In degrees $1$ through $8$ its only nonzero homogeneous pieces have degrees $4,6,8$; they are one-dimensional, generated by $F,H,F^2$, respectively. Each $S_k|_W$ is invariant. It therefore remains only to determine three constants.

For clarity, the calculation can be performed by the following exact recurrence, with coefficients in $\mathbb Z[s]$ and $s^2=-7$. Put
$P_{0,\varnothing}(t)=1$, and, for $|A|=r$, put
\[
 P_{r,A}(t)=\sum_{j\in A}
       \bigl(1+t(B_0)_{r-1,j}\bigr)P_{r-1,A\setminus\{j\}}(t).
\]
Expanding along the last row proves that
$P_{8,P}(t)=\operatorname{perm}(J+tB_0)$. Dividing its coefficient of $t^k$ by $(8-k)!$ gives
\[
\begin{array}{c|l}
 k&S_k(B_0)\\\hline
 4&83328+8064s\\
 6&-2502656+2396160s\\
 8&18391040+3809280s.
\end{array}
\]
Since $F(1,0,0)=1$ and $H(1,0,0)=-216(3+s)$, these values give the stated $c,d$ and $5/1764$. The same values follow from the inclusion-exclusion formula for the permanent.

Consequently a nonzero $B(x)$ is a level direction precisely when
$F(x)=H(x)=0$. The quartic is smooth, and its intersection with its Hessian consists of its twenty-four ordinary flexes
\cite[\S2.1, Proposition (i), pp.~64--65]{Elkies1999}; see also
\cite[Lemma 1.4, p.~32]{Kaneta1998}.
A level plane would give a projective line contained in this smooth irreducible quartic, which is impossible. Conjugation gives the other summand.
\end{proof}

The invariant ideal of these level conditions has a particularly short description. With Klein's generators $\Phi_4,\Phi_6,\Phi_{14},\Phi_{21}$,
\[
 I=(\Phi_4,\Phi_6)\subset\mathbb C[W]^G,
 \qquad
 \mathbb C[W]^G/I\simeq
 \mathbb C[\Phi_{14},\Phi_{21}]/(\Phi_{21}^2-\Phi_{14}^3).
\]
The last relation is Elkies's equation (1.18) modulo $(\Phi_4,\Phi_6)$.
Thus the invariant ideal is prime, although its extension to the polynomial ring cuts out the cone over twenty-four distinct points.

\begin{proposition}\label{prop:klein-mixed}
For $T=B(x)+\overline B(y)$, where $x,y\in\mathbb C^3$ are independent coordinates, put $q(x,y)=\sum_i x_i y_i$. Then
\[
 S_2(T)=1792q(x,y),
\]
\[
 S_4(T)=cF(x)+\overline c\,\overline F(y)+774144q(x,y)^2.
\]
Every level subspace of $W\oplus\overline W$ has dimension at most two.
\end{proposition}

\begin{proof}
The quadratic formula follows by polarizing
$S_2(T)=\frac12\sum_{i,j}T_{ij}^2$ on doubly centred matrices and using the basis pairings above. For the quartic, the representation identities
$\operatorname{Sym}^2W\simeq\mathbf6$ and
$\operatorname{Sym}^3W\simeq\mathbf7\oplus W^*$
\cite[p.~55]{Elkies1999}
show that the invariant spaces in bidegrees $(3,1)$ and $(1,3)$ vanish, while the invariant space in bidegree $(2,2)$ is spanned by $q^2$. The coefficient of $q^2$ is determined by
\[
 S_4(B_0+\overline B_0)=940800,
 \qquad c+\overline c=166656;
\]
their difference is $774144$. These evaluations use the preceding recurrence.

The quadratic condition first gives $\dim L\leq3$. Suppose equality holds.
Neither projection of $L$ to a pure summand can have rank at most one, because its kernel would contain a two-dimensional pure level subspace, excluded by Theorem~\ref{thm:klein-pure}.
If the projection to $W$ has rank three, then $L$ is the graph of $y=Ax$.
Isotropy gives $A^{\mathsf T}=-A$, hence $\operatorname{rank}A\leq2$.
The quartic identity becomes $cF(x)+\overline c\,\overline F(Ax)=0$.
It would make $F$ constant along a nonzero direction in $\ker A$.
A homogeneous plane quartic independent of one coordinate is singular at the corresponding point, contrary to the smoothness of $F=0$.
The same argument applies with the two projections reversed.

Both projections therefore have rank two. Their one-dimensional kernels in $L$ are distinct. The equality of the two restricted quartics makes either polynomial constant along both kernel directions. It is consequently a nonzero fourth power of a linear form on $L$. Nonzero follows because the Klein quartic contains no line. Hence the projective line in either projection meets the quartic at a single point with multiplicity four. This would be a hyperflex. Klein's twenty-four flexes are all ordinary, so it has no hyperflex. This contradiction proves the bound.
\end{proof}

\begin{proposition}\label{prop:klein-noextension}
If a level subspace of $W\oplus\overline W$ contains a nonzero pure level direction from either summand, then it has dimension one.
\end{proposition}

\begin{proof}
It suffices to exclude a level plane $L$ containing a line $\ell\subset W$.
The projection to $\overline W$ has rank one. If the other projection had rank two, the quartic identity would express the restriction of $F$ to a projective line as either zero or a fourth power, again giving a line on the quartic or a hyperflex. Both are impossible. Thus
$L=\ell\oplus\ell'$ for a pure level line $\ell'\subset\overline W$, and the two lines must be orthogonal for $q$.

There are only twenty-four possible $\ell'$. The following finite calculation excludes them without approximating their coordinates. Let $\zeta$ be a primitive seventh root of unity, use
$s=\zeta+\zeta^2+\zeta^4-\zeta^3-\zeta^5-\zeta^6$, and reduce at the prime given by $\zeta\mapsto16$ in $\mathbb F_{29}$; then $s\mapsto15$.
By transitivity we may take $\ell$ to be the $\zeta$-eigenline of $\rho(u)$.
Its reduction is $x=(1,16,3)$. Generate the conjugate orbit using the two displayed generator matrices with $s$ replaced by $-s$, starting at the $\zeta^{-1}$-eigenline. There are twenty-four distinct reductions. Precisely two have dot product zero with $x$, and the permanent recurrence gives
\[
\begin{array}{c|c}
 y&S_6\bigl(aB(x)+b\overline B(y)\bigr)\pmod {29}\\\hline
 (1,15,26)&8a^4b^2\\
 (1,19,24)&11a^2b^4.
\end{array}
\]
Every generator matrix is integral at this prime and invertible after reduction. The orbit has twenty-four points in characteristic zero, and its twenty-four distinct reductions therefore represent every point injectively. An orthogonality relation in characteristic zero survives reduction, so the two displayed candidates include every possible partner. Their nonzero sextic reductions rule out a vanishing sextic in characteristic zero. This excludes the plane; conjugation handles the other summand. The exact orbit enumeration and permanent computations are reproduced by the programs in the workbench.
\end{proof}

A mixed level plane transverse to both pure summands is not excluded by these arguments. Its existence is a separate finite-dimensional problem.

\paragraph{Other finite-group slices.}
Let $G\leq S_n\times S_n$ be finite and let $W$ be an isotypic summand for an irreducible representation not isomorphic to its dual. There are no invariant linear or quadratic forms on $W$, hence $S_1|_W=S_2|_W=0$. Put
$R=\mathbb C[W]^G$ and
$I=(S_3|_W,\ldots,S_n|_W)\subset R$.
Then
\[
 I\mathbb C[W]\cap R=I.
\]
Indeed, averaging an expression $f=\sum h_kS_k$ for an invariant $f$ replaces each $h_k$ by its Reynolds average and preserves the equality. Thus the quotient ideal is exactly the contraction of the level ideal. Each degree $k$ with $R_k=0$ contributes no equation; an explicit invariant presentation reduces the remaining equations to finitely many homogeneous components. The Klein calculation above carries this procedure through for one conjugate pair, without asserting a classification for arbitrary finite groups.
